\section{\sysname{} Design}

We design the \sysmodlong{} (\sysmod{}), based on the definitions in previous section.
A \sysmod{} represents a single distributed system with a specific semantics
and \sysname{} can perform individual \sysmod{} validation to check if all
semantic constraints are sound and compatible, as
well as pairwise checks between two \sysmods{} to determine if one
is satisfiable under another's constraints. 
More importantly, \sysname{} can model applications that rely on multiple 
distributed systems with different
consistency semantics (i.e., \sysmods{}) and can check for compatibility 
between the application logic and the assigned semantics by performing a 
compositional check using the \sysgraph{}.

\subsection{Compatibility Check}

A key feature of \sysname{} is the ability to reason not only about
compatibility between consistency semantics. 
We define compatibility between two consistency semantics using \sysmods{} 
 $M_1$ and $M_2$ using
an implication-based criterion. $M_1$ is considered compatible with
$M_2$ if the formula $M_1 \Rightarrow M_2$ is valid, i.e., there is no
execution that satisfies all constraints of $M_1$ without also
satisfying $M_2$.
We implement this by asserting the negation $\neg (M_1 \Rightarrow M_2)$ and
checking if it is unsatisfiable using our SMT-based approach. If no
counterexample can be found, it implies that $M_1$
refines/subsumes $M_2$. This compatibility check is not symmetric:
$M_1 \Rightarrow M_2$
holding does not necessarily mean that $M_2 \Rightarrow M_1$ also holds.

Figure \needref{} shows a result of pairwise compatibility checks for
base semantics in \sysname{} which in-part shows that our consistency
models compatibility faithfully matches with the theoretical results
in \cref{fig:semantics-tree}.

To combine multiple consistency models into a single stronger model, we used
logical conjunctions on each \sysmod{}'s constraints. By taking the constraints
representing each model's semantics and forming their conjunction, we derive a
composed model that enforces all included constraints simultaneously. 
This approach allows incremental composition, users can start from a
base model and iteratively strengthen it by adding new sets of constraints
(either from our implemented models or user-defined constraints) representing
additional consistency guarantees. 


\subsection{\sysgraph{} and Composition Check}

%. Components (e.g., functionality) 
%of an application are modeled as nodes and data flow and constraints 
%between systems are represented as edges connecting the nodes in the 
%\sysgraph{}. \sysname{} can check specific assignments of distributed
%system semantics or outputs results of all satisfiable assignments 
%for the \sysgraph{}.

Applications that rely on multiple distributed systems with different
consistency semantics can be checked for correctness by performing a 
compositional check using the \sysgraph{}.
Our composability check
verifies if a distributed application, modeled as a graph (\ie 
\sysgraph{}) can have its components interact in a way that
respects their individual consistency requirements and guarantees.

%The function checks if there exists at least one valid way to assign
%specific semantics (from the choices provided in needs/provs/cons)
%such that all interactions are consistent. 

The \sysgraph{} is composed of \textit{nodes} with two fields: 
(1) \textit{needs}, required constraint tuples from dependent components, and 
(2) \textit{provs}, constraint tuples it provides to other nodes. 
While user provided application logics can have both needs and provs, the 
built-in \sysmods{} only specify provs as their consistency semantics.
\textit{Edges} connect nodes and have the \textit{cons} field, 
specific constraint tuples for the node pair interaction. 
The tuples within the list represent constraints that must hold together 
(logical and). 

In the implementation, we perform a Depth-First Search (DFS) starting from a
given source node, keeping track of visited edges and the accumulated
constraints (\texttt{path\_premise}) along the current path. 
\sysname{} performs compatible checks, \texttt{compatible(check\_provs,
compose(check\_needs, check\_ec), path\_premise)}:
for each edge (u,v) from the current node u, it checks if provs from v
and cons from the edge (u, v) satisfy needs from u. 
This checks if the chosen path provides guarantee of the destination (check\_provs) is
implied by the conjunction (compose) of the chosen needs of the
source (check\_needs) and the chosen edge constraints (check\_ec),
considering the constraints already accumulated (path\_premise). 

If all check succeeds then the application is proven to be compatible
with the plugged in \sysmods{} and the application logics are also
sound as a whole. 

\subsection{Compatible Combination Search}

\sysname{} is capable of searching for compatible \sysmods{} for an
application given a \sysgraph{}. When blank \sysmods{} are plugged into the 
\sysgraph{}, \sysname{} searches for compatible combinations of 
\sysmods{} for the user.
Considering that strictly weaker distributed semantics can be replaced
with strictly stronger ones and the weaker semantics are typically 
more performant and cheaper to maintain, \sysname{} searches the space
in an order that prioritize weaker semantics. Specifically, we give scores
1 to 6 to semantics in the columns from left to right in \cref{fig:semantics-tree}
and test the combination that has lower sum of scores first during the search. 
If a compatible combinations are found, the search can stop early; replacing a 
\sysmod{} in the compatible combination with a strictly stronger one is always 
guaranteed to be compatible. 
However, our search strategy is not optimal and investigating more 
efficient search algorithm remains as future work. 



%If a compatible combination is found that satisfies all unvisited edges
%between u and v at the current step, it adds the chosen constraints
%(needs, provs, cons) to the result graph, updates the path\_premise
%by composing the newly satisfied constraints, marks the edges as
%visited, and continues the DFS from the destination node v. If the
%DFS eventually visits all edges in the graph, the function returns
%True and the result graph showing one valid assignment. If a dead end
%is reached or no compatible combination works for an edge, it
%backtracks, removing the choices made at that step and trying other
%combinations. If the DFS completes without visiting all edges
%(meaning no fully consistent assignment was found), it returns False.

%% TODO: can be better
%After a successful composable check, our extraction function
%(\texttt{extract(inode, onode, result)}) takes the resulting graph
%from composability check (which represents one valid assignment of
%constraints) and aggregates a single Z3 constraint that
%represents the net or equivalent semantic guarantee of the subgraph
%reachable from a specific node (inode), assuming inode and onode are
%the same. The function performs a DFS (dfs helper) starting from the
%specified inode on this result graph. For each edge (src, dst)
%traversed in the DFS, the extraction function uses the specific provs
%constraint assigned to src in the result graph, uses the specific
%needs constraint assigned to dst in the result graph, and retrieves
%the specific cons constraint assigned to the edge (src, dst, k) in
%the result graph. The extract function creates a composite constraint
%for this single edge interaction: compose(src\_provs, edge\_cons,
%dst\_needs). This represents the conjunction of constraints relevant
%to this specific step in the flow. It collects these composite
%constraints from all edges reachable from inode in the DFS. Finally,
%it returns the logical and (compose) of all collected composite
%constraints. Simply put, axioms and relations define constraints on
%Operations; similarly, extract defines a summary constraint on a
%successfully composed subgraph. If the composable check found a valid
%way for components A, B, and C to interact starting from A,
%extract(A, A, result) generates a single formula representing the
%overall guarantee provided by the A→B→C... chain, as seen from A's
%perspective, based on the specific choices made during the composability check.
%
%
